\begin{result}{E.2.31}{An explicit cusp conjugator outside the proper completion}{Proved}
Retain the actual groups \(F=\langle x,y\rangle\),
\(G=\widehat F\), and \(\mathcal H_c\subset G\) of E.2.30.
There is a specified element \(f\in G\) such that
\[
 f^{-1}yf\in\mathcal H_c,
 \qquad
 f\notin\mathcal H_c C_G(y).
\tag{E.2.31a}
\]
Its specification below uses the actual \(A_5\)-cover of
E.2.29, an explicit invariant proper homology class, and
integral transfer along a fixed cofinal cover tower. Thus
membership of one conjugated cusp generator does not force
its conjugator to lie in \(\mathcal H_c\) times the cusp
centralizer. No assertion is made that this \(f\) defines an
automorphism preserving all three peripheral classes or
occurs as the conjugator of an arithmetic Galois automorphism.
\end{result}
\begin{proof}
We construct the compatible element using integral transfer
on the original compactified covers, rather than choosing
unrelated local conjugators.

\paragraph{1. Actual fixed spaces and integral transfer.}
Use the exact groups and extensions from E.2.30:
\[
 1\longrightarrow\widehat A_N^c\longrightarrow P_N^c
   \longrightarrow Q_N\longrightarrow1,
 \qquad A_N^c=H_1(C_N,\Z),\qquad\Gamma_N^c=F/L_N.
\]
Let \(y_N\) be the actual image of \(y\) in \(P_N^c\), and set
\[
 Z_N^y=(\widehat A_N^c)^{y_N},\qquad
 T_N^y=\Gamma_N^c Z_N^y\subset P_N^c.
\tag{E.2.31b}
\]
The latter is a subset; no subgroup assertion is needed.
Conjugation by \(y_N\) on the abelian kernel is its actual
deck action. For \(\gamma\in\Gamma_N^c\) and \(z\in Z_N^y\),
\[
 (\gamma z)y_N(\gamma z)^{-1}
      =\gamma y_N\gamma^{-1}\in\Gamma_N^c.
\tag{E.2.31c}
\]

For \(M\subset N\), the actual compactified map
\(\pi:C_M\to C_N\) has degree \(e=[N:M]\). It may be
branched at the filled cusps. Its integral pushforward on
\(H_1\) is exactly the restriction of the quotient map
\(P_M^c\to P_N^c\) to the abelian kernels: both are induced
by the original map of covers and their filling maps.
Define the integral umkehr using oriented Poincar\'{e} duality:
\[
 \pi^!=\operatorname{PD}_{C_M}^{-1}\pi^*
              \operatorname{PD}_{C_N}:
 H_1(C_N,\Z)\longrightarrow H_1(C_M,\Z).
\tag{E.2.31d}
\]
Then \(\pi_*\pi^!=e\,\mathrm{id}\). Indeed, pairing
\(\pi_*\pi^!a\) with any integral class \(b\) equals the
cup-product pairing of \(\pi^*\operatorname{PD}(a)\) and
\(\pi^*\operatorname{PD}(b)\) on \(C_M\). This is \(e\)
times their pairing on \(C_N\), by the oriented degree
formula. Unimodularity of the integral intersection pairing
proves the stated integral identity, also for a branched map.

The actual deck maps satisfy \(\pi y_M=y_N\pi\).
Naturality of oriented duality and pullback therefore gives
\[
 \pi^!(A_N^c)^{y_N}\subset(A_M^c)^{y_M},
 \qquad\pi_*\pi^!=e\,\mathrm{id}.
\tag{E.2.31e}
\]
The invariant lattice is the kernel of the integer matrix
\(y_N-1\). Smith normal form, using injectivity of every
nonzero integer on \(\widehat\Z\), proves
\[
 Z_N^y=(A_N^c)^{y_N}\otimes\widehat\Z.
\tag{E.2.31f}
\]
After completion the actual image of the fixed-space
transition thus satisfies
\[
 \pi_*(Z_M^y)\supset eZ_N^y.
\tag{E.2.31g}
\]

\paragraph{2. Surjectivity of the compatible sets.}
The maps \(T_M^y\to T_N^y\) are onto. For
\(\gamma z\in T_N^y\), choose an integral invariant
\(a_0\in(A_N^c)^{y_N}\) representing \(z\) modulo
\(eZ_N^y\); this exists by (E.2.31f). Then
\[
 z'=\pi^!((z-a_0)/e)\in Z_M^y,
 \qquad\pi_*z'=z-a_0.
\]
Lift \(\gamma a_0\) through the surjection
\(\Gamma_M^c\to\Gamma_N^c\), obtaining \(\gamma'\).
The product \(\gamma' z'\) maps exactly to \(\gamma z\).
This proves surjectivity of these sets even though the
fixed-space transition itself need not be surjective.

\paragraph{3. The actual starting cover and fixed integral class.}
Use the monodromy of E.2.29, with \(x=(123)\),
\(y=(12345)\), and the same permutation-composition
convention as its frozen certificate. These generate
\(A_5\), and \(xy\) has order five. The regular cover
has degree sixty and inertia orders \((3,5,5)\), so
\[
 2g(C)-2
  =60\bigl(-2+(1-1/3)+(1-1/5)+(1-1/5)\bigr)=16.
\]
Thus \(g(C)=9\). Let \(N_0=\ker(G\to A_5)\) and
\(H=\langle y\rangle\), of order five.

Each order-five branch fibre consists of twelve cosets
\(A_5/I\). Exactly two are fixed by \(H\): the cosets
with \(gIg^{-1}=H\) are counted by
\(\lvert N_{A_5}(H)\rvert/\lvert I\rvert=10/5=2\).
All order-five subgroups are conjugate. There are no
\(H\)-fixed points in the order-three fibre. Each
nonidentity element of \(H\) generates \(H\), hence
fixes exactly four points of \(C\), two over each
order-five branch point. Riemann--Hurwitz for the actual
quotient gives
\[
 16=5\bigl(2g(C/H)-2\bigr)+4(5-1),
 \qquad g(C/H)=1.
\tag{E.2.31h}
\]
In particular the rational invariant proper homology has
dimension two. The following explicit integral calculation
also supplies its required nonzero class directly.

In the actual sixty-one-dimensional affine cycle basis of
E.2.29, let \(B\) be the saturated boundary lattice, let
\(Y\) be the deck matrix of \(y\), and let \(E\) be its
integer averaged boundary projector, so \(E^2=60E\).
Put \(Q=60I-E\), which kills \(B\), and define
\[
 R=(I+Y+Y^2+Y^3+Y^4)Q.
\tag{E.2.31i}
\]
The first nonzero column of \(R\), with zero-based index
one, is the vector \(w\) specified by the frozen matrix
data and the standalone certificate. The exact checks are
\[
 \operatorname{rank}R=2,\quad Yw=w,\quad Ew=0,\quad
 w_0=9,\qquad qB=0,\quad qw=540,
 \quad q=Q_{0,*}.
\tag{E.2.31j}
\]
Therefore the proper class \(v=[w]\in A_{N_0}^c\) is
nonzero, integral, and \(y\)-invariant. The displayed
integer functional \(q\) descends to the actual proper
lattice and detects \(v\) by \(540\).

Choose the explicit \(a\in\widehat\Z\) with
\(a_2=0\) and \(a_p=1\) for every odd prime, and set
\[
 z_0=av\in Z_{N_0}^y.
\tag{E.2.31k}
\]
It is not integral: an integral value of its detector
would be the integer \(540a\). Its \(2\)-adic coordinate
is zero, forcing that integer to be zero, while its
\(3\)-adic coordinate is the nonzero element \(540\).
Since \(\Gamma_{N_0}^c\cap\widehat A_{N_0}^c=A_{N_0}^c\),
one also has \(z_0\notin\Gamma_{N_0}^c\).

\paragraph{4. A completely specified compatible element.}
Order reduced words in the alphabet
\(x,x^{-1},y,y^{-1}\) by length and then lexicographically
in the displayed alphabet order. Starting with \(N_0\),
let \(N_n\) be its intersection with all kernels of
homomorphisms \(G\to S_j\) for \(j\leq n\). This
specified descending normal chain is cofinal; repeated
levels are harmless.

Put \(g_0=1\), and start with \(z_0\) from (E.2.31k).
Suppose \(g_n\in F\) and \(z_n\in Z_{N_n}^y\) are
defined. Write \(e_n=[N_n:N_{n+1}]\), and choose the
first word \(u_n\in K_{N_n}\) whose proper class
\(a_n\) is \(y\)-invariant and satisfies
\(z_n-a_n\in e_n Z_{N_n}^y\). Such a word exists by
(E.2.31f) and the surjection
\(K_{N_n}\to A_{N_n}^c\). Define
\[
 \begin{aligned}
 g_{n+1}&=g_nu_n,\\
 z_{n+1}&=\pi_n^!((z_n-a_n)/e_n),\\
 t_n&=(g_n\bmod L_{N_n})z_n\in T_{N_n}^y.
 \end{aligned}
\tag{E.2.31l}
\]
The image of \(u_n\) at level \(N_n\) is \(a_n\),
and the image of \(z_{n+1}\) is \(z_n-a_n\).
Consequently \(t_{n+1}\) maps exactly to \(t_n\).
The actual identification
\(G\simeq\varprojlim P_{N_n}^c\) of E.2.30 supplies
one unique \(f_0\in G\) with these coordinates.
No computability assertion about arbitrary profinite
coordinates is part of this specification.

The initial coordinate is \(z_0\), so
\(f_0\notin\mathcal H_c\). Formula (E.2.31c)
places every coordinate of \(f_0yf_0^{-1}\) in
\(\Gamma_{N_n}^c\). Cofinality gives
\[
 f_0\notin\mathcal H_c,
 \qquad f_0yf_0^{-1}\in\mathcal H_c.
\tag{E.2.31m}
\]

\paragraph{5. The exact centralizer and the failed implication.}
For a specified free generator of a free profinite group,
the original Corollary 2.3 of
\cite[p.~332]{GenBuxErshovRapinchuk2011} proves
\[
 C_G(y)=\overline{\langle y\rangle}=y^{\widehat\Z}.
\tag{E.2.31n}
\]
Its proof uses the explicit finite relation-module
Proposition 2.2(i), pp.~330--331, and applies to this
actual generator and completion. By E.2.30,
\(y^{\widehat\Z}\subset\mathcal H_c\), so
\(\mathcal H_c C_G(y)=\mathcal H_c\).
Setting \(f=f_0^{-1}\) in (E.2.31m) proves (E.2.31a).

The finite certificate is
\path{work/code_generation/check_cusp_conjugator_seed.py}.
It checks the frozen E.2.29 JSON hash, actual group and
normalizer counts, all displayed genus identities, and
the integral matrix identities in (E.2.31j). Its input
JSON has SHA-256
\texttt{ba0e10ed1fc1eb306fafdfbb2e77fe3ed38f17f5ba0f5593188f59203872b111}.
The infinite element and its all-level properties are
proved by (E.2.31l), not inferred from the finite check.
The assignments \(x\mapsto x\), \(y\mapsto f^{-1}yf\)
still require an actual automorphism, the third peripheral
relation, and arithmetic Galois compatibility before they
can address the full E.2.22 gate.
\end{proof}
